\documentclass[11pt]{article}

\usepackage[a4paper, margin=1.5cm]{geometry}
\usepackage{amsmath}
\usepackage{amssymb}
\usepackage{multicol}
\usepackage{booktabs}
\usepackage{enumitem}
\usepackage[colorlinks=true, linkcolor=blue, urlcolor=blue]{hyperref}
\usepackage{xcolor}
\usepackage{tikz}

\setlength{\parindent}{0pt}
\setlength{\parskip}{4pt}
\pagestyle{empty}

\definecolor{accent}{RGB}{40,80,160}
\definecolor{lightgray}{RGB}{240,240,240}

\newcommand{\sechead}[1]{\vspace{6pt}{\color{accent}\textbf{\large #1}}\vspace{2pt}\hrule\vspace{4pt}}
\newcommand{\Mcl}{\mathcal{M}_{\mathrm{cl}}}
\newcommand{\Sadm}{\mathcal{S}_{\mathrm{adm}}}
\newcommand{\phig}{\varphi}

\begin{document}

\begin{center}
{\LARGE\bfseries Constraint-Manifold Physics (VFD)}\\[4pt]
{\large\color{accent} A structural framework for mass, confinement, and geometry}\\[2pt]
{\small Institute of Vibrational Field Dynamics\quad|\quad vibrationalfielddynamics.org\quad|\quad April 2026}
\end{center}

\vspace{4pt}
\hrule height 1.5pt
\vspace{8pt}

\begin{multicols}{2}

\sechead{Core Idea}

Instead of starting from fields and dynamics, this framework starts from:
\begin{itemize}[nosep, leftmargin=12pt]
    \item a fixed geometric structure (600-cell, $\phig$-scaled)
    \item explicit closure constraints on admissible states
\end{itemize}

\textbf{Physical states are exactly those configurations that satisfy closure.}\\
Everything else is excluded.

\sechead{Key Equation}

\begin{equation*}
\boxed{\;m = \kappa \cdot d(\psi,\, \Mcl)\;}
\end{equation*}

\begin{itemize}[nosep, leftmargin=12pt]
    \item $\Mcl$: constraint manifold (closed states)
    \item $d(\psi, \Mcl)$: distance to closure
    \item Mass = distance from constraint satisfaction
\end{itemize}

\sechead{What It Explains (Structurally)}

\begin{itemize}[nosep, leftmargin=12pt]
    \item \textbf{Mass} --- closure residual (zero fitted parameters, 13 particles at 0.014\% avg)
    \item \textbf{Confinement} --- only connected states are admissible
    \item \textbf{Electroweak} --- projection invariance (torsional symmetry analogue)
    \item \textbf{Geometry} --- closed states form a manifold
    \item \textbf{Gravity} (analogue) --- curvature of that manifold
    \item \textbf{Fine structure constant} --- from 600-cell eigenvalues (0.81 ppm)
\end{itemize}

\sechead{The Progression}

\begin{tabular}{@{}rl@{}}
Paper V & particle masses from spectral geometry \\
Paper VI & electroweak as boundary projection \\
Paper VII & closure operator formalised \\
Paper VIII & confinement as connectivity \\
Paper IX & continuous geometry \\
Paper X & gravitational analogy \\
Paper XI & Standard Model correspondence \\
\end{tabular}

\sechead{Why This Might Matter}

\begin{itemize}[nosep, leftmargin=12pt]
    \item Reduces 13 mass parameters to zero fitted
    \item Unifies mass, confinement, gauge, gravity under one system
    \item Provides an alternative geometric starting point
    \item Is \textbf{falsifiable} under improved precision
\end{itemize}

\columnbreak

\sechead{Standard Model Mapping}

{\small
\begin{tabular}{@{}ll@{}}
\toprule
\textbf{Standard Model} & \textbf{Constraint Framework} \\
\midrule
Hilbert space & $\Sadm$ \\
Physical states & $\Mcl$ (constraint manifold) \\
Gauge symmetry & projection invariance \\
Mass (Higgs) & $m = \kappa\, d(\psi, \Mcl)$ \\
Massless & on $\Mcl$ ($\Delta = 0$) \\
Massive & off $\Mcl$ ($\Delta > 0$) \\
Confinement & connectivity constraint \\
Colour (3) & threefold shell structure (analogue) \\
Spacetime (GR) & $\Mcl$ (analogue only) \\
Geodesics & constrained motion on $\Mcl$ \\
Gravity & manifold curvature \\
\midrule
Lagrangian & \textit{not provided} \\
Renormalisation & \textit{not addressed} \\
Scattering & \textit{not computable yet} \\
\bottomrule
\end{tabular}
}

\sechead{What This Is NOT}

\begin{itemize}[nosep, leftmargin=12pt]
    \item[$\times$] Not a derivation of the Standard Model
    \item[$\times$] Not a replacement for QFT
    \item[$\times$] Not a derivation of general relativity
    \item[$\times$] Not a complete physical theory (no dynamics yet)
\end{itemize}

This is a \textbf{structural framework}, not a finished theory.\\
\textit{Stage: pre-dynamical (constraint-level only).}

\sechead{How to Test / Falsify}

\begin{itemize}[nosep, leftmargin=12pt]
    \item Mass correspondences break with improved precision
    \item Constraint filtering fails to restrict states correctly
    \item Discrete $\to$ continuous mapping becomes inconsistent
    \item Constraint system fails to produce stable composite classes
\end{itemize}

\sechead{Entry Points}

\textbf{Start here:} Paper XI (SM correspondence)\\
\textbf{Mass result:} Paper V (13 particles, 0.014\%)\\
\textbf{Bridge:} Paper IV (proton-electron, $\phig^{1265/81}$)

\vspace{4pt}
{\small
\url{https://github.com/vfd-org/vfd-crystallisation}\\
\url{https://vibrationalfielddynamics.org}
}

\end{multicols}

\vspace{4pt}
\hrule height 1.5pt
\vspace{4pt}
\begin{center}
{\small\itshape ``If this is wrong, it should break under tighter constraints.''}
\end{center}

\end{document}
